\section{Scope and method}
\label{sec:method}

\paragraph{Scope.}
\todo{What is included and what is not, and why.}

\paragraph{Finding the sources.}
\todo{Where the sources came from (databases, citation chasing, the project literature catalogue),
the search terms, and the date of the search. Say honestly what was not searched.}

\paragraph{Comparing the methods.}
Every method is described by the same columns, in the style of the summary table of
Singh et~al.~\cite{singh2019}:

\begin{itemize}
  \item \textbf{Measurements used}: the sensor signals the method needs.
  \item \textbf{Estimated quantity}: what the method produces or decides.
  \item \textbf{Model used}: the vehicle, tyre or process model it assumes, if any.
  \item \textbf{Methodology}: the estimation or control technique.
  \item \textbf{Reference}: where to read about it.
\end{itemize}

\noindent Columns can be added when the topic needs them (for instance validation level:
simulation, offline data or on-vehicle).
